\chapter{Conclusions, Limitations and Future Work}\label{ch:conclusions}

\needswork{STATUS: DRAFTED. The future-work section is the one most likely to be examined for
evidence that the candidate understands the work's boundaries; it is written accordingly and
deliberately does not oversell.}

\section{What this thesis establishes}

\begin{enumerate}
\item After compression, authentication --- not data --- is the dominant on-air cost of a
tamper-evident UAV telemetry ledger. At the smallest encoding studied, the signature is larger than
the record it signs.

\item Frame-level self-batching is the correct placement, and it Pareto-dominates block-level
aggregation under \emph{any} stationary loss process, not merely under the independent-loss
assumption the emulator implements.

\item On IEEE 802.11 the freshness deadline binds before the frame size does, so the amplification
law that motivates aggressive compression is never operative: compression pays exactly $1\times$.
On low-rate links the frame size binds and the leverage is real. \textbf{The regimes are different
in kind.}

\item Three of seven EU868 data rates admit no per-frame-verifiable telemetry at any encoding,
batch, header or chain design, because a $64$\,B signature does not fit a $51$\,B payload. This is
arithmetic and cannot move.

\item Joined to a validated channel model, the co-designed configuration sustains
$1.9$--$3.2\times$ the neighbourhood of an inline baseline, at $58.7\%$ fewer on-air bytes, at an
operating point compliant with 3GPP TS~22.125.

\item Of the four design axes, only placement and batching genuinely interact. The four-axis
co-design claim, as originally stated, was overstated.
\end{enumerate}

\section{Limitations, stated as such}

\subsection{Modelling scope}

The single collision domain admits no spatial reuse and no hidden terminals; a real formation has
both, one helping and one hurting. Mobility was tested and found null for collision-limited
capacity, but per-frame Doppler is unmodelled --- the null result means mobility does not change
\emph{collision-limited capacity}, not that mobility is harmless to a link. Loss is independent per
frame; the ordering result survives correlation but absolute verifiability figures do not. The
keyframe interval is fixed and unoptimised. Post-quantum sizes are projections, not
implementations.

\subsection{Validation scope}

Simulation validates \emph{throughput}. Latency, energy and loss behaviour are not validated by it,
and ``simulation-validated'' must not be read as more. Hardware validates \emph{airtime} on a
two-node link; contention remains simulation-only. The utilisation ceiling is measured at one node
count. Composed energy is a lower bound by roughly $10$--$14\%$, and the residual is uncharged
prototype framing overhead --- measured, stated, not corrected away.

\subsection{The low-rate arm}

Analysis and simulation only: no radio metered, no energy column, one spreading factor per run.
The capacity figure is a within-one-spreading-factor bound, not a network capacity, and its
confidence interval is a knife edge.

\subsection{The boundary itself}

Three of the four excluded data rates are arithmetic. The fourth was excluded by this design's own
frame header, and this thesis says so and quantifies the recovery. A reader should treat the
three-rate claim as durable and the fourth as an engineering statement about \emph{this} wire
format.

\section{Future work}

\subsection{Immediately available, and cheap}

\begin{description}
\item[Invariance of the utilisation ceiling in $N$.] Frame-size invariance was measured; node-count
invariance was not. One additional sweep closes it, and the scenario already exists.

\item[Pricing implicit certificates.] The certificate accounting charges an explicit certificate
periodically. Implicit (ECQV) certificates are the standard smaller alternative and were never
modelled. The current charge is an upper bound, which cuts against this work's own comparison ---
so the gap is safe but should be closed.

\item[An integer-keyed wire profile, adopted rather than measured.] \Cref{ch:lowrate} measures it
and declines to adopt it, because the format is frozen for reproducibility. A successor should
adopt it, re-freeze, and report the boundary that results --- which will be a \emph{smaller}
exclusion and a better-engineered system.
\end{description}

\subsection{Substantial, and worth doing}

\begin{description}
\item[Channel-access delay in the freshness model.] The freshness constraint that \emph{sets} the
batch omits the dominant delay term at high load. It is currently mitigated by reporting
utilisation alongside, not by fixing the model. A DCF access-delay term would make the freshness
constraint credible at high utilisation, where it currently is not.

\item[Multi-spreading-factor capacity on the low-rate arm.] Each run fixes one data rate,
forfeiting the quasi-orthogonality that supplies much of a real deployment's capacity. That answers
a different question --- deployment capacity rather than per-configuration feasibility --- and it
is the question an operator would ask next.

\item[Contention on real hardware.] Two radios validate airtime; the contention model, which
carries the capacity envelope, remains simulation-only. A rig with enough nodes to reach the
interesting region of the curve is the single most valuable missing measurement in this thesis.

\item[A post-quantum arm that is implemented rather than projected.] The projections show that
batching cannot rescue post-quantum signature sizes, because freshness caps the batch and one
signature plus header already exceeds the frame budget. That is a strong negative result and it
deserves measurement rather than arithmetic.
\end{description}

\subsection{Speculative, and stated as such}

The exclusion bound is link-agnostic. Instantiating it on other constrained bands --- and on the
signature sizes the post-quantum standards actually mandate --- would say which of tomorrow's radio
and cryptographic combinations are feasible before either is deployed. That is the natural
generalisation, and it is the one this author would pursue.

\section{Closing}

This thesis introduces no new cryptographic primitive and makes no claim to one. Its contribution
is a reproducible co-design of existing, standardised components; the quantified trade-offs that
co-design entails; the boundary beyond which no choice of those components helps at all; and an
account --- unusually explicit for a thesis --- of the errors made in establishing that boundary and
the classes of defect that survived the safeguards intended to prevent them.

The most durable result here is also the smallest one that survived scrutiny: on the longest-range
EU868 data rates, a $64$-byte signature does not fit a $51$-byte payload, and nothing in the design
space changes that.
